\documentclass[11pt]{article}

\usepackage[margin=1in]{geometry}
\usepackage{amsmath,amssymb}
\usepackage{booktabs}
\usepackage{array}
\usepackage{longtable}
\usepackage{hyperref}
\usepackage{setspace}
\setstretch{1.08}

\title{Nimbus Advanced Agricultural Benchmarking}
\author{Sheerwater Benchmark · Method Note\\August 2026}
\date{}

\begin{document}
\maketitle
\thispagestyle{empty}

\paragraph{TL;DR.}
\textbf{Actionable days beyond baseline}:
for each grid cell, take the longest lead at which it still meets the
actionable percent-good cutoff, subtract the same horizon for the baseline,
then average over all cells in the selected country or countries.

\section{Motivation}

In the Nimbus program, our north star is that smallholder farmers across
sub-Saharan Africa have the forecasts they need, backed by their national
meteorological and hydrological services, to make critical agricultural
decisions during the rainy season. Benchmarking must help us ask whether
forecasts are \textbf{good enough} for those decisions: scores should look good
when a farmer could have used the forecast well, and bad when it would have
misled them.

Standard daily rainfall metrics (MAE, CRPS, brittle contingency scores) often
fail that test: they mix wet and dry regimes, punish tiny timing shifts with
double penalties, or hinge on hard thresholds that do not match farmer-relevant
``wet enough'' / ``dry enough'' questions. Sheerwater's agricultural event
metrics instead verify specific decision-relevant events with temporal
softening so that rain a day or two off is not scored as a full miss.

\subsection*{Agriculturally relevant events}

In our first implementation, we focus on the following events. ``Good enough''
is defined separately for each event (magnitude, dryness, timing, or
accumulation), and can be updated as application needs change.

{\small
\begin{longtable}{@{}>{\raggedright\arraybackslash}p{0.16\textwidth}
                  >{\raggedright\arraybackslash}p{0.20\textwidth}
                  >{\raggedright\arraybackslash}p{0.28\textwidth}
                  >{\raggedright\arraybackslash}p{0.26\textwidth}@{}}
\toprule
\textbf{Event} & \textbf{Agricultural decision} & \textbf{Definition}
  & \textbf{How can we measure good enough?} \\
\midrule
\endfirsthead
\toprule
\textbf{Event} & \textbf{Agricultural decision} & \textbf{Definition}
  & \textbf{How can we measure good enough?} \\
\midrule
\endhead
Extreme rain days
  & Flood / erosion risk; whether to delay field work or protect harvested crops
  & Days where 1-day rainfall is at or above the local 90th percentile of
    daily precip from IMERG Final, 1998--2025
  & Score = SMAPE of rainfall in the 5-day window centered on the event
    (\(\pm 2\) days). \textbf{Good enough:} SMAPE \(\le 0.30\)
    (within about 30\% of the true rainfall in that window) \\
\addlinespace
In-season dry spells
  & When to plant; whether to irrigate, replant, or expect moisture stress during crop establishment
  & After the season has started (first \(>4\)\,mm rain day near 10\% of
    seasonal accumulation) and before 40\% of seasonal rain has fallen
    (seasons with \(<150\)\,mm total rain excluded): any 10-day window with
    \(<15\)\,mm total rain
  & Score = forecast rainfall (mm) in that same 10-day window.
    \textbf{Good enough:} forecast rain \(\le 20\)\,mm
    (at most 20\,mm predicted where observations were dry) \\
\addlinespace
Early-season rainfall accumulation
  & When to plant; whether early-season moisture is on track for crop establishment
  & First day in the season when cumulative rainfall reaches \(100\)\,mm;
    the evaluated quantity is rainfall accumulated over the preceding
    30 days (right-aligned). For forecast leads shorter than 30 days,
    earlier days are filled from observations
  & Score = MAE (mm) of that 30-day accumulation vs truth.
    \textbf{Good enough:} MAE \(\le 30\)\,mm
    (the forecasted 30-day total is within 30\,mm of the observed total) \\
\addlinespace
Start of season date\(^*\)
  & When to plant / prepare land for the rainy season
  & First date in each season meeting a start-of-season rule
    (ICPAC, CHC, or Moron \& Robertson)
  & Score = absolute error (days) between forecast and observed
    start-of-season date. Good-enough bar TBD \\
\bottomrule
\end{longtable}
}

\vspace{-0.5em}
{\footnotesize \(^*\)Implementation in progress.}

\paragraph{Reading the good-enough bars.}
For extreme rain days we use
\(\mathrm{SMAPE}=|F-O|/(|F|+|O|)\).
SMAPE \(\le 0.30\) means the absolute error is at most 30\% of the combined
forecast and observed magnitude. In plain terms, the forecast rainfall in the
softened window was within about 30\% of the true rainfall.
For dry spells, \(\le 20\)\,mm means the forecast did not put more than 20\,mm
of rain into a 10-day window that was observed dry (\(<15\)\,mm).
For early-season accumulation, MAE \(\le 30\)\,mm means the forecast's 30-day
total missed the observed total by at most 30\,mm.

\subsection*{Observation-triggered and forecast-triggered events}

For event metrics such as extreme rain days and dry spells, an event can be
defined from the observations or from the forecast. We score both:

\begin{itemize}
  \item \textbf{Observation-triggered:} events found in the observations.
        Poor scores here are \emph{false negatives}: rain (or a dry spell)
        happened, but the forecast missed it.
  \item \textbf{Forecast-triggered:} events found in the forecast.
        Poor scores here are \emph{false positives}: the forecast warned of
        an event that was not observed.
\end{itemize}

KR Tracking's default ``both'' mode combines the two by taking the
\textbf{worse} of the observation-triggered and forecast-triggered scores
at each cell and lead (equivalently: the maximum error, or the minimum
percent-good). That way a forecast cannot look good by catching observed
events while over-warning, or by staying quiet while missing real events;
both failure modes count.

\vspace{0.5em}

The \textbf{KR Tracking} dashboard then summarizes those
event-based scores as:
\begin{enumerate}
  \item how many grid cells clear a chosen quality cutoff at each lead
        (\textbf{good cells at a cutoff}), and
  \item how far each cell's skill extends in lead time relative to a baseline
        (\textbf{days beyond baseline}).
\end{enumerate}
This note defines those two quantities.

\section{Setup}

Let \(F\) be a candidate forecast and \(B\) a baseline forecast
(e.g.\ climatology or a reference NWP product). For a chosen agricultural
event metric, each of \(F\) and \(B\), at lead day \(L\) and grid cell, has a
\textbf{percent-good} score in \([0,1]\): the fraction of verified events
meeting the metric's good-enough bar.

Fix two cutoffs on that score.
These cutoffs are flexible and can be changed for different applications or
risk tolerances; for now we use \(\theta_{\mathrm{info}}=0.50\) and
\(\theta_{\mathrm{act}}=0.75\):

\begin{center}
\begin{tabular}{@{}llp{0.52\textwidth}@{}}
\toprule
Cutoff & Symbol & Meaning \\
\midrule
Informative & \(\theta_{\mathrm{info}}=0.50\)
  & At least 50\% of events at that cell and lead are well predicted.
    Useful signal (worth watching), but usually not enough to drive
    operational forecast alerts or firm farm decisions on its own. \\
\addlinespace
Actionable & \(\theta_{\mathrm{act}}=0.75\)
  & At least 75\% of events are well predicted.
    High enough consistency to support decision-grade use
    (e.g.\ alerts, irrigation timing, or planting guidance). \\
\bottomrule
\end{tabular}
\end{center}

Leads considered for the horizon KPI are
\[
L \in \{7,14,21,28\},
\]
i.e.\ forecast weeks 1--4. We do not evaluate weeks 5 and 6 (\(L=35,42\))
here, because skill at those leads was consistently very low.
A cell \textbf{passes} cutoff \(\theta\) at lead \(L\) when its percent-good is
at least \(\theta\).

\subsection*{Probabilistic forecasts}

Ensemble forecasts use the same event definitions and good-enough bars as
deterministic forecasts. We do \textbf{not} collapse the ensemble to a single
mean field before scoring. Instead, each ensemble member is scored with the
deterministic metric and marked good or bad against the good-enough bar; the
resulting good-enough percentage is then averaged across all ensemble members.
KR Tracking can be run in deterministic or probabilistic mode (separate cached
tables). In probabilistic mode, cell curves and days-beyond KPIs use this
member-averaged rate: how often a draw from the ensemble would have been good
enough, not whether the ensemble-mean rainfall field alone cleared the bar.

\section{Good cells at a cutoff}

At a fixed lead \(L\), define
\[
G_F(L,\theta)
  \;=\;
  \#\{\text{cells with percent-good}\ge\theta\}.
\]
The dashboard's threshold curve plots \(G_F(L,\theta)\) versus \(\theta\).
The Informative and Actionable markers are narrow guides at
\(\theta_{\mathrm{info}}\) and \(\theta_{\mathrm{act}}\).
The lead-panel deltas are
\[
\Delta G(L,\theta)
  \;=\;
  G_F(L,\theta) - G_B(L,\theta),
\]
reported as a \textbf{number of cells} versus baseline \(B\).

Optional table views also report \(G_F(L,\theta_{\mathrm{info}})\) and
\(G_F(L,\theta_{\mathrm{act}})\) directly
(``Cells @ Informative / Actionable'').
A separate ``Expected Cells'' view is an AUC-like summary of the threshold
curve: the average of \(G_F(L,\theta)\) over all cutoffs \(\theta \in [0,1]\).
It is a curve-wide score, not the days-beyond KPI.

\section{Days beyond baseline}

Days beyond baseline answers: \emph{relative to the baseline, how much farther
in lead time does this forecast stay good enough?}

\subsection*{Per cell}

For each grid cell \(c\) and cutoff
\(\theta \in \{\theta_{\mathrm{info}},\theta_{\mathrm{act}}\}\),
define the \textbf{horizon} as the longest lead at which the cell still passes:
\begin{equation}
H_F(c,\theta)
  \;=\;
  \max\bigl\{\,
    L \in \{7,14,21,28\}
    :\;
    \text{percent-good}_F(c,L) \ge \theta
  \,\bigr\},
\end{equation}
or \(0\) if the cell never clears \(\theta\) at those leads.
A null percent-good at a lead is a noop: that lead is dropped for both
\(F\) and \(B\) on the cell (horizons use only shared usable leads). It is
not treated as a fail.

The cell's days beyond baseline are
\begin{equation}
D_F(c,\theta)
  \;=\;
  H_F(c,\theta) - H_B(c,\theta),
\end{equation}
with both horizons computed on the intersection of leads where \(F\) and
\(B\) have usable percent-good. Cells with no shared usable leads are
omitted from the mean.

\subsection*{Region summary}

Regions are selectable as a single country or multiple countries. The reported
value is the \textbf{mean over all grid cells} in the selected country or
countries:
\begin{equation}
\mathrm{Days}(F,B;\theta)
  \;=\;
  \frac{1}{n}\sum_{c} D_F(c,\theta),
\end{equation}
where the sum runs over the \(n\) comparable cells in that selection.
Separate values are shown for Informative (\(\theta_{\mathrm{info}}\)) and
Actionable (\(\theta_{\mathrm{act}}\)).
The baseline column is always \(0\).
Forecasts with no overlapping usable cells show as missing, not as zero
(so short-range or sparse models are not scored as ``tied with baseline'').

Against a baseline with horizon 0 everywhere, that mean is easy to read:
every cell actionable out to lead 28 gives \(+28\); every cell actionable only
at lead 7 gives \(+7\); if 3 of 15 cells are actionable out to lead 14 and the
rest are unchanged, the region mean is \(3\times 14/15=+2.8\); and matching
the baseline everywhere gives \(0\).
A cell-level \(+14\) means that cell clears the cutoff two weeks farther than
the baseline on that same cell.

\section{Limitations and next steps}

\begin{itemize}
  \item \textbf{Training / testing conflict.}
        Several ML forecasts were trained on years that overlap the
        evaluation period (e.g.\ training through 2021), so KR Tracking
        scores can be optimistic relative to a true out-of-sample holdout.
        \emph{Planned:} a rolling overfit sensitivity study that varies the
        evaluation window relative to each model's training cutoff.
  \item \textbf{Probabilistic evaluation.}
        Ensemble forecasts are scored per member and then averaged, which
        captures \textbf{expected} error across the ensemble rather than the
        true error of a single realized forecast or a proper score of the
        full distribution.
        \emph{Planned:} compute a median event from the ensemble, then evaluate
        that deterministic event.
  \item \textbf{Equal cell weighting.}
        Regional means treat every grid cell equally, so empty or sparsely
        farmed cells weigh as much as intensive agricultural areas.
        \emph{Planned:} weight cells by arable land or farming population.
  \item \textbf{Discrete weekly leads.}
        Horizons are evaluated only at four leads
        (\(L \in \{7,14,21,28\}\): weeks 1--4). Sub-weekly skill is ignored:
        a cell that remains good enough out to day 10 is still tracked as
        horizon 7.
        \emph{Planned:} denser lead sampling (e.g.\ daily or every few days)
        so mid-week gains are counted.
\end{itemize}

\vspace{1em}
\noindent Implemented in the
\href{https://dashboards.rhizaresearch.org/d/ee4mze492j0n4d/home?orgId=1\&from=now-6h\&to=now\&timezone=utc}{KR Tracking dashboard}.

\end{document}
